% ------------------------------------------------------------------------
%
\begin{frame}{Term rewriting}

  More general are the \textbf{term\hyp{}rewriting systems}, where a
  \textbf{term} is a mathematical object built on tuples, integers and
  variables.

  \bigskip

   Let us consider the following totally ordered system:
   \begin{align*}
     (0,m) &\xrightarrow{1} m\\
     (n,m) &\xrightarrow{\smash 2} (n-1, n \times m)\\
     n     &\xrightarrow{\smash 3} (n,1)
   \end{align*}

   Arithmetic operators \((-)\)~and~\((\cdot)\) are defined out of the
   system and \(m\)~and~\(n\) are variables denoting natural numbers.

   \bigskip

   Would the rules not be ordered as they are actually laid out, the
   second rule would match any pair. Instead, it can be assumed that
   \(n \neq 0\) when matching with it.

   \bigskip

   The last rule is to be tried last, as it matches all terms.

\end{frame}

% ------------------------------------------------------------------------
%
\begin{frame}{Transitive closures}

  We can easily see that all the compositions of rewrites starting
  with a natural number~\(n\) end with the factorial of~\(n\):
  \smallskip
  \begin{equation*}
    n \xrightarrow{3} (n,1)
    \xrightarrow{2} \dots
    \xrightarrow{2} (0,n!)
    \xrightarrow{1}
    n!, \quad \text{for \(n \in \mathbb{N}\)}.
  \end{equation*}

  Let us note \((\xrightarrow{n})\) the composition of
  \((\rightarrow)\) repeated \(n-1\)~times:
  \begin{align*}
    (\xrightarrow{1})   &:= (\rightarrow)\\
    (\xrightarrow{\smash{n+1}}) & :=
    (\rightarrow) \circ (\xrightarrow{\smash{n}}),
    \quad\text{with \(n > 0\)}.
  \end{align*}

  The \textbf{transitive closure} of \((\rightarrow)\) is defined as
  \begin{equation*}
    (\twoheadrightarrow) := \bigcup_{i > 0}{(\xrightarrow{i})}.
  \end{equation*}

  The factorial function coincides with the transitive closure of
  \((\rightarrow)\): \(n \twoheadrightarrow n!\). Let
  \((\xrightarrow{*})\) be the reflexive\hyp{}transitive closure of
  \((\rightarrow)\), that is,
  \begin{equation*}
    (\xrightarrow{*}) := (=) \cup (\twoheadrightarrow).
  \end{equation*}

\end{frame}

% ------------------------------------------------------------------------
%
\begin{frame}{Functions}

  \begin{itemize}

    \item A confluent system defines a \textbf{partial function}.

    \item If the system is also terminating, we get a
      \textbf{function}.

    \item We can name functions: for example, \(\fun{c}(1,
      \fun{d}(n))\)~is a term constructed with \textbf{function names}
      \fun{c}~and~\fun{d}, as well as variable~\(n\).

    \item A tuple tagged with a function name, like \(\fun{f}(x,y)\),
      is called a \textbf{function call}.

    \item The components of the tuples are then called
      \textbf{arguments}, for example \(\fun{d}(n)\) is the second
      argument of the call \(\fun{c}(1, \fun{d}(n))\).

    \item It is possible for a function call to hold no arguments,
      like \(\fun{d}()\).

    \item We restrict the left\hyp{}hand sides of rules to be function
      calls.

  \end{itemize}

\end{frame}

% ------------------------------------------------------------------------
%
\begin{frame}{Trees}

  The topological understanding of a function call or a tuple is the
  finite \textbf{tree}. A tree is a hierarchical layout of
  information, for example:
  \begin{center}
    \includegraphics{tree_for_term}
  \end{center}

   The disks are called \textbf{nodes} and the segments which connect
   two nodes are called \textbf{edges}. The topmost node (with a
   diameter) is called the \textbf{root} and the bottommost ones
   (\(\bullet\)) are called the \textbf{leaves}. All nodes except the
   leaves are seen to downwardly connect to some other nodes, called
   \textbf{children}. Upwardly, each node but the root is connected to
   another node, called its \textbf{parent}.

\end{frame}

% ------------------------------------------------------------------------
%
\begin{frame}{Trees}

  \textbf{Trees are pervasive} in informatics, especially in
  functional programming, but not only. Structured documents, like
  books, articles, and web pages, are composed of chapters, sections,
  paragraphs, figures, appendices, indices, etc.

  \bigskip

  The occurrences of these components are mutually constrained; for
  instance, it is understood that a section is part of a chapter and
  that appendices are located at the end of a document.

  \bigskip

  This hierarchical layout is meant to facilitate reading, and it
  supports the search for specific items of information.

  \bigskip

  When considering computer systems, these data must be uniformly
  encoded by means of a formal language.

\end{frame}

% ------------------------------------------------------------------------
%
\begin{frame}{Terms as trees}

  \label{terms_as_trees}

  Trees can be used to depict terms as follows. A function call is a
  tree whose root is the function name and the children are the trees
  denoting the arguments.

  \bigskip

  A tuple can be considered as having an invisible function name
  represented by a node with a period (\texttt{.}) in the tree, in
  which case the components of the tuple are its children. For
  example,
  \begin{center}
    \includegraphics{tree}
  \end{center}
  is the tree corresponding to the tern
  \(\fun{f}((\fun{g}(0),(x,1)),(),\fun{g}(y))\).

  \bigskip

  Note that tuples, except the empty tuple, are not represented in the
  tree, since they encode the structure itself, which is already laid
  out. The number of arguments of a function is called \textbf{arity}.

\end{frame}

% ------------------------------------------------------------------------
%
\begin{frame}{Functional languages}

  \begin{itemize}

    \item For programming, we only consider confluent systems because
      they define partial functions.

    \item To have only one reduction chain per term to the normal
      form, we may also enforce that rules are ordered.

    \item We also enforce that normal forms must be \textbf{values},
      that is, terms without function calls or only with calls to
      functions not occurring on the left\hyp{}hand sides.

    \item In other words, irreducible calls of defined functions are
      not allowed in values.

    \item Systems with these constraints are called \textbf{functional
      rewrite systems}.

    \item A program here is a term, and evaluating it means to rewrite
      the term
      \begin{enumerate}
        \item to a value,
        \item or to end in error with a normal form that is not a value,
        \item or to never terminate.
      \end{enumerate}

  \end{itemize}

\end{frame}

% ------------------------------------------------------------------------
%
\begin{frame}{Call-by-value}

  If a function call contains at least one other function call, there
  is a choice as to which one to reduce first.

  \bigskip

  One reduction strategy, named \textbf{call-by-value}, consists in
  always reducing the arguments before the call itself. Unfortunately,
  call-by-value enables otherwise terminating rewrites to not
  terminate. For instance, let us consider
  \begin{equation*}
    \fun{f}(x) \xrightarrow{\alpha} 0\qquad
    \fun{g}() \xrightarrow{\beta} \fun{g}()
  \end{equation*}
  We have \(\fun{f}(\fun{g}()) \xrightarrow{\alpha} 0\) but
  \(\fun{f}(\fun{g}()) \xrightarrow{\beta} \fun{f}(\fun{g}())
  \xrightarrow{\beta} \dots\)

  \bigskip

  Despite the loss of expressiveness, we shall retain in this lecture
  the call by value, as it is the strategy of \OCaml and it is easy to
  follow.

  \bigskip

  In general, the order in which the arguments of a function call are
  evaluated is not specified.

\end{frame}

% ------------------------------------------------------------------------
%
\begin{frame}{Call-by-value}

  Another reason to choose call-by-value is that it allows us to
  restrict the shape of the left\hyp{}hand sides, called
  \textbf{patterns}, to one, outermost function call.

  \bigskip

  For instance, we can then disallow, for being useless, a
  rule like

  \begin{equation*}
    \fun{plus}(x,\fun{plus}(y,z)) \rightarrow
    \fun{plus}(\fun{plus}(x,y),z).
  \end{equation*}

  \noindent because, following call\hyp{}by\hyp{}value, the argument
  \(\fun{plus}(y,z)\) must be evaluated before the outermost call,
  therefore it cannot be part of any pattern.

\end{frame}

% ------------------------------------------------------------------------
%
\begin{frame}{Termination}

  \label{termination2}

  We do not require by construction that all functional systems
  terminate, although that is a desirable property, but we
  nevertheless speak of `functions' instead of `partial functions'
  (which would be technically correct).

  \bigskip

  If we enforce termination, we loose
  \textbf{Turing\hyp{}completeness}: we cannot program all computable
  functions.

  \bigskip

  The termination of rewrite systems and, in particular, functional
  programs, is an \textbf{undecidable problem} in general.

  \bigskip

  Nevertheless, there exists a set of well\hyp{}known \textbf{decision
    criteria}, which we will use when programming recursive functions
  in \OCaml.

\end{frame}

% ------------------------------------------------------------------------
%
\begin{frame}{Higher-order programs, recursion, $\lambda$-calculus}

  Most functional languages allow \textbf{higher\hyp{}order function}
  definitions, whereas standard term\hyp{}rewriting systems do
  not. For example:
  \begin{equation*}
    \fun{f}(g,0) \rightarrow 1
    \qquad
    \fun{f}(g,n) \rightarrow n \times g(g,n-1)
    \qquad
    \fun{fact}(n) \rightarrow \fun{f}(\fun{f},n)
  \end{equation*}
  where \(n \in \mathbb{N}\). Note that these two definitions are
  \textbf{not} recursive, yet the function \fun{fact} is the factorial
  function. A function definition is \textbf{recursive} if the name of
  the definition occurs on at least one right\hyp{}hand side.

  \bigskip

  An adequate theoretical framework to understand higher\hyp{}order
  functions is \textbf{\(\lambda\)-calculus}. In fact,
  \(\lambda\)-calculus features prominently in the semantics of
  programming languages, even not functional ones.

  \bigskip

  For didactic purposes, we began with rewrite systems because they
  offer implicit pattern matching (that is, rule selection).

\end{frame}
